% GENERATED FILE. Edit canonical lessons, then run npm run generate:books.
% canonical-source-hash: fnv1a64:ce1595499c05d211
% canonical-lessons: RU-C253-skazka, RU-C253-istoriya, RU-C253-shakhmaty, RU-C253-transport, RU-C253-trolleybus

\chapter{Fairy tale, Story, Chess, Transport, Trolleybus}
\label{ch:ru-fairy-tale-story-chess-transport-trolleybus}
\hlchaptermodality{253}

\begin{chapteropening}
\textbf{By the end of this chapter:} \emph{I can say a fairy tale, a story, chess, transport and a trolleybus.}

Complete the last lesson of chapter 253: I can say a fairy tale, a story, chess, transport and a trolleybus.
\end{chapteropening}

\section[\texorpdfstring{\ru{сказка} (skázka)}{сказка (skázka)}]{\ru{сказка} (skázka) — a fairy tale}

\label{lesson:RU-C253-skazka}



\begin{quote}
Before the new one: say the Russian for running, then the Russian for swimming.
\end{quote}

\subsection*{You'll want to know: \ru{сказка}}
\textbf{\ru{сказка}} (\emph{skázka}) — \textquotedblleft{}a fairy tale\textquotedblright{}.

\begin{cousinweb}
\textbf{\ru{сказка}} comes from \emph{*sŭkazŭka}, from \emph{\ru{казать}}, \textquotedblleft{}to show, to say\textquotedblright{} (Vasmer); it is recorded from the 1600s, before which a tale was a \emph{\ru{баснь}}.
\end{cousinweb}

\begin{grammarlens}[title={What you've built}]
One more word for this chapter.
\end{grammarlens}

\subsection*{Your turn}
\begin{itemize}
\raggedright
  \item \emph{Say it:} \emph{skázka}
  \item \emph{Say it:} \emph{skázka}, once more
  \item \emph{Say it:} say \emph{skázka}
  \item \emph{Recall:} say the Russian for running, then the Russian for swimming, then say \emph{skázka} again
\end{itemize}

\begin{tcolorbox}[breakable,colback=teal!4,colframe=teal!35!black,title={Before you move on}]
What does \ru{сказка} mean? (\textquotedblleft{}A fairy tale\textquotedblright{}.) Say it once more.
\end{tcolorbox}
\section[\texorpdfstring{\ru{история} (istóriya)}{история (istóriya)}]{\ru{история} (istóriya) — a story; history}

\label{lesson:RU-C253-istoriya}



\begin{quote}
Before the new one: say the Russian for swimming, then the Russian for a fairy tale.
\end{quote}

\subsection*{You'll want to know: \ru{история}}
\textbf{\ru{история}} (\emph{istóriya}) — \textquotedblleft{}a story; history\textquotedblright{}.

\begin{cousinweb}
\textbf{\ru{история}} — from Peter the Great's time; perhaps through German \emph{Historie} from Latin \emph{historia}, Greek \emph{historía}, \textquotedblleft{}inquiry, a narrative\textquotedblright{} (Vasmer).
\end{cousinweb}

\begin{grammarlens}[title={What you've built}]
One more word for this chapter.
\end{grammarlens}

\subsection*{Your turn}
\begin{itemize}
\raggedright
  \item \emph{Say it:} \emph{istóriya}
  \item \emph{Say it:} \emph{istóriya}, once more
  \item \emph{Say it:} say \emph{istóriya}
  \item \emph{Recall:} say the Russian for swimming, then the Russian for a fairy tale, then say \emph{istóriya} again
\end{itemize}

\begin{tcolorbox}[breakable,colback=teal!4,colframe=teal!35!black,title={Before you move on}]
What does \ru{история} mean? (\textquotedblleft{}A story; history\textquotedblright{}.) Say it once more.
\end{tcolorbox}
\section[\texorpdfstring{\ru{шахматы} (shákhmaty)}{шахматы (shákhmaty)}]{\ru{шахматы} (shákhmaty) — chess}

\label{lesson:RU-C253-shakhmaty}



\begin{quote}
Before the new one: say the Russian for a fairy tale, then the Russian for a story; history.
\end{quote}

\subsection*{You'll want to know: \ru{шахматы}}
\textbf{\ru{шахматы}} (\emph{shákhmaty}) — \textquotedblleft{}chess\textquotedblright{}.

\begin{cousinweb}
\textbf{\ru{шахматы}} — Old Russian \emph{\ru{шахматы}}, from 1280; it is built on \emph{\ru{шах} \ru{мат}}, \textquotedblleft{}checkmate\textquotedblright{}, from Persian \emph{shāh}, \textquotedblleft{}king\textquotedblright{}, probably through a German source (Vasmer).
\end{cousinweb}

\begin{grammarlens}[title={What you've built}]
One more word for this chapter.
\end{grammarlens}

\subsection*{Your turn}
\begin{itemize}
\raggedright
  \item \emph{Say it:} \emph{shákhmaty}
  \item \emph{Say it:} \emph{shákhmaty}, once more
  \item \emph{Say it:} say \emph{shákhmaty}
  \item \emph{Recall:} say the Russian for a fairy tale, then the Russian for a story; history, then say \emph{shákhmaty} again
\end{itemize}

\begin{tcolorbox}[breakable,colback=teal!4,colframe=teal!35!black,title={Before you move on}]
What does \ru{шахматы} mean? (\textquotedblleft{}Chess\textquotedblright{}.) Say it once more.
\end{tcolorbox}
\section[\texorpdfstring{\ru{транспорт} (tránsport)}{транспорт (tránsport)}]{\ru{транспорт} (tránsport) — transport}

\label{lesson:RU-C253-transport}



\begin{quote}
Before the new one: say the Russian for a story; history, then the Russian for chess.
\end{quote}

\subsection*{You'll want to know: \ru{транспорт}}
\textbf{\ru{транспорт}} (\emph{tránsport}) — \textquotedblleft{}transport\textquotedblright{}.

\begin{cousinweb}
\textbf{\ru{транспорт}} — from 1720; it came from French or Dutch \emph{transport} or German \emph{Transport} (Vasmer), from Latin \emph{trānsportāre}, \textquotedblleft{}to carry across\textquotedblright{}.
\end{cousinweb}

\begin{grammarlens}[title={What you've built}]
One more word for this chapter.
\end{grammarlens}

\subsection*{Your turn}
\begin{itemize}
\raggedright
  \item \emph{Say it:} \emph{tránsport}
  \item \emph{Say it:} \emph{tránsport}, once more
  \item \emph{Say it:} say \emph{tránsport}
  \item \emph{Recall:} say the Russian for a story; history, then the Russian for chess, then say \emph{tránsport} again
\end{itemize}

\begin{tcolorbox}[breakable,colback=teal!4,colframe=teal!35!black,title={Before you move on}]
What does \ru{транспорт} mean? (\textquotedblleft{}Transport\textquotedblright{}.) Say it once more.
\end{tcolorbox}
\section[\texorpdfstring{\ru{троллейбус} (trolléybus)}{троллейбус (trolléybus)}]{\ru{троллейбус} (trolléybus) — a trolleybus}

\label{lesson:RU-C253-trolleybus}



\begin{quote}
Before the new one: say the Russian for chess, then the Russian for transport.
\end{quote}

\subsection*{You'll want to know: \ru{троллейбус}}
\textbf{\ru{троллейбус}} (\emph{trolléybus}) — \textquotedblleft{}a trolleybus\textquotedblright{}.

\begin{cousinweb}
\textbf{\ru{троллейбус}} comes from English \emph{trolleybus}: \emph{trolley}, the wheel on a pole that runs along the wire, with \emph{bus}.
\end{cousinweb}

\begin{grammarlens}[title={What you've built}]
One more word for this chapter.
\end{grammarlens}

\subsection*{Your turn}
\begin{itemize}
\raggedright
  \item \emph{Say it:} \emph{trolléybus}
  \item \emph{Say it:} \emph{trolléybus}, once more
  \item \emph{Say it:} say \emph{trolléybus}
  \item \emph{Recall:} say the Russian for chess, then the Russian for transport, then say \emph{trolléybus} again
\end{itemize}

\begin{tcolorbox}[breakable,colback=teal!4,colframe=teal!35!black,title={Before you move on}]
What does \ru{троллейбус} mean? (\textquotedblleft{}A trolleybus\textquotedblright{}.) Say it once more.
\end{tcolorbox}
